% BHU PYQ -- AIHC & Archaeology: Political History of India (Circa 600 BCE - 600 CE) (2024-25 NEP)
\documentclass[11pt,a4paper]{article}
\usepackage[a4paper,top=2.5cm,bottom=2.5cm,left=2.8cm,right=2.8cm,headheight=14pt]{geometry}
\usepackage{amsmath,amssymb}
\usepackage{fontspec}
\setmainfont{Kohinoor Devanagari}
\usepackage{microtype,fancyhdr,enumitem,hyperref,lastpage,parskip}

\hypersetup{colorlinks=true,urlcolor=black,linkcolor=black,pdftitle={AIHMJ101: Political History of India (600 BCE - 600 CE) -- BHU NEP 2024-25}}

\pagestyle{fancy}\fancyhf{}
\renewcommand{\headrulewidth}{0.4pt}\renewcommand{\footrulewidth}{0.4pt}
\fancyhead[L]{\small \textit{Banaras Hindu University}}
\fancyhead[R]{\small \textit{AIHMJ101: Political History of India (2024-25)}}
\fancyfoot[C]{\small Page \thepage\ of \pageref{LastPage}}

\newlist{parts}{enumerate}{2}
\setlist[parts,1]{label=(\alph*),leftmargin=*,itemsep=4pt,topsep=3pt}
\setlist[parts,2]{label=(\roman*),leftmargin=*,itemsep=2pt,topsep=2pt}

\newcommand{\pts}[1]{\hfill \textit{[#1]}}

\begin{document}

\begin{center}
  {\large\bfseries BANARAS HINDU UNIVERSITY}\\[2pt]
  {\normalsize B. A. (Hons.) Semester I Examination, 2024-25 (NEP)}\\[6pt]
  \rule{0.8\linewidth}{0.5pt}\\[6pt]
  {\large\bfseries ANCIENT INDIAN HISTORY, CULTURE \& ARCHAEOLOGY}\\[4pt]
  {\large\bfseries Paper Code: AIHMJ-101 : Political History of India (Circa 600 BCE -- 600 CE)}\\[4pt]
  {\normalsize Time: 2:30 Hours \hfill Full Marks: 70}\\[2pt]
  {\small REF NO. 102}
\end{center}
\rule{\linewidth}{0.4pt}
\smallskip

\noindent \textit{Instructions: Attempt all questions. Questions of Section A, B and C carry 10, 05 and 02 marks each respectively. Write your Roll No.\ at the top immediately on receipt of this question paper.}

\smallskip
\noindent \textit{सभी प्रश्नों के उत्तर दीजिए। खण्ड अ, ब तथा स के प्रश्न क्रमशः 10, 05 एवं 02 अंकों के हैं।}

\bigskip

\section*{SECTION -- A / खण्ड -- अ}
\noindent \textit{Note: Long Answer Type Questions. Answer each in approximately 250 words.}\pts{3 $\times$ 10 = 30}

\noindent \textit{दीर्घ उत्तरीय प्रश्न। प्रत्येक प्रश्न का उत्तर लगभग 250 शब्दों में दीजिए।}

\begin{enumerate}[label=\textbf{\arabic*.},leftmargin=*,itemsep=12pt]

  \item Describe the political condition of the Sixteen Mahajanapadas of the Sixth Century BCE.\pts{10}
  
  छठी शताब्दी ई० पू० के सोलह महाजनपदों की राजनीतिक स्थिति का वर्णन कीजिए।

  \begin{center}\textbf{OR / अथवा}\end{center}

  Give a detailed account of Alexander's Invasion on India.\pts{10}
  
  भारत पर सिकन्दर के आक्रमण का विस्तृत विवरण दीजिए।

  \item Evaluate the contribution of the Haryanka dynasty to the expansion of the Magadha Empire.\pts{10}
  
  मगध साम्राज्य के विस्तार में हर्यंक वंश के योगदान का मूल्यांकन कीजिए।

  \begin{center}\textbf{OR / अथवा}\end{center}

  Discuss the political and cultural achievements of the Gupta ruler, Chandragupta II.\pts{10}
  
  गुप्त शासक `चंद्रगुप्त द्वितीय' की राजनीतिक एवं सांस्कृतिक उपलब्धियों की विवेचना कीजिए।

  \item Throw light on Ashoka's Dhamma.\pts{10}
  
  अशोक के धम्म पर प्रकाश डालिए।

  \begin{center}\textbf{OR / अथवा}\end{center}

  Who were the Kushanas? Determine the extent of the Kushana empire during the reign of Kanishka I.\pts{10}
  
  कुषाण कौन थे? कनिष्क प्रथम के शासनकाल में कुषाण साम्राज्य की सीमाओं का निर्धारण कीजिए।

\end{enumerate}

\bigskip

\section*{SECTION -- B / खण्ड -- ब}
\noindent \textit{Note: Short Answer Type Questions. Answer each in approximately 125--150 words.}\pts{4 $\times$ 5 = 20}

\noindent \textit{लघु उत्तरीय प्रश्न। प्रत्येक प्रश्न का उत्तर लगभग 125--150 शब्दों में दीजिए।}

\begin{enumerate}[label=\textbf{\arabic*.},leftmargin=*,start=4,itemsep=10pt]

  \item Describe the achievements of Pushyamitra Shunga.\pts{5}
  
  पुष्यमित्र शुंग की उपलब्धियों का वर्णन कीजिए।

  \begin{center}\textbf{OR / अथवा}\end{center}

  Write a short note on Indo-Greek Kings.\pts{5}
  
  हिन्द-यवन (इण्डो-ग्रीक) राजाओं पर एक संक्षिप्त टिप्पणी लिखिए।

  \item Critically analyse Samudragupta's policy towards the conquered states.\pts{5}
  
  विजित राज्यों के प्रति समुद्रगुप्त की नीति का समालोचनात्मक विश्लेषण कीजिए।

  \begin{center}\textbf{OR / अथवा}\end{center}

  Discuss the career and achievements of Rudradaman.\pts{5}
  
  रुद्रदामन के जीवन एवं उपलब्धियों की चर्चा कीजिए।

  \item Sketch the career and achievements of Kharavela on the basis of Hathigumpha Inscription.\pts{5}
  
  हाथीगुम्फा अभिलेख के आधार पर खारवेल के जीवन एवं उपलब्धियों का चित्रण कीजिए।

  \begin{center}\textbf{OR / अथवा}\end{center}

  Discuss the achievements of Chandragupta Maurya.\pts{5}
  
  चंद्रगुप्त मौर्य की उपलब्धियों की विवेचना कीजिए।

  \item Write a note on Persian Invasion on India.\pts{5}
  
  भारत पर पारसी (ईरानी) आक्रमणों पर एक टिप्पणी लिखिए।

  \begin{center}\textbf{OR / अथवा}\end{center}

  Describe the causes of the downfall of the Mauryan dynasty.\pts{5}
  
  मौर्य साम्राज्य के पतन के कारणों का वर्णन कीजिए।

\end{enumerate}

\bigskip

\section*{SECTION -- C / खण्ड -- स}
\noindent \textit{Note: Very Short Answer Type Questions. Answer each in 15--20 words.}\pts{10 $\times$ 2 = 20}

\noindent \textit{अति लघु उत्तरीय प्रश्न। प्रत्येक प्रश्न का उत्तर 15--20 शब्दों में दीजिए।}

\begin{enumerate}[label=\textbf{\arabic*.},leftmargin=*,start=8,itemsep=8pt]
  \item 
  \begin{parts}
    \item Name the inscriptions that describe the Sudarshan Lake.\pts{2}
    
    सुदर्शन झील का वर्णन करने वाले अभिलेखों के नाम लिखिये। (Junagadh Rock Inscription of Rudradaman / Skandagupta)

    \item Who was Mihirakula?\pts{2}
    
    मिहिरकुल कौन था? (Huna ruler)

    \item What is the importance of the ``Bhitari Inscription''?\pts{2}
    
    `भीतरी स्तम्भ अभिलेख' का क्या महत्त्व है? (Skandagupta's victory over Hunas)

    \item Who was Heliodorus?\pts{2}
    
    हेलियोडोरस कौन था? (Greek ambassador who erected Garuda Pillar at Besnagar)

    \item Locate Shravasti.\pts{2}
    
    श्रावस्ती की स्थिति बताइये। (Capital of Kosha Mahajanapada, Uttar Pradesh)

    \item What is the policy of ``Grahanmokshanugraha''?\pts{2}
    
    ``ग्रहणमोक्षानुग्रह'' की नीति क्या है? (Samudragupta's southern campaign policy)

    \item What is Sallekhana?\pts{2}
    
    सल्लेखना (संथारा) क्या है? (Jain ritual of fasting to death)

    \item Name two books written by Vishakhadatta.\pts{2}
    
    विशाखदत्त की दो पुस्तकों के नाम लिखिये। (Mudrarakshasa, Devichandraguptam)

    \item Who was Ramagupta?\pts{2}
    
    रामगुप्त कौन था? (Gupta ruler mentioned in Devichandraguptam)

    \item What do you understand by the term ``Apasirve''?\pts{2}
    
    ``अपसिन्वे'' (Apasinve) शब्द से आप क्या समझते हैं?
  \end{parts}
\end{enumerate}

\end{document}
